% Hand-wrapped from mineru OCR of pages 199-200 (Theorem 3.6 + proof)
% Source: Shao, Mathematical Statistics (linear regression chapter)
% Correction kind: env_wrap (raw OCR uses "Theorem 3.6." plain-text header)
\documentclass{article}
\usepackage{amsmath,amssymb,amsthm}
\newtheorem{theorem}{Theorem}
\begin{document}

% --- Context (model 3.25) ---
% Linear model: $X = Z \beta + \varepsilon$, where $Z$ is an $n \times p$ design
% matrix of rank $r \le p$. There exists an $n \times r$ submatrix $Z_*$ of $Z$
% of rank $r$ and a fixed $r \times p$ matrix $Q$ such that
% \[ Z = Z_* Q. \tag{3.30} \]
% An LSE of $\beta$ has the form $\hat{\beta} = (Z^\tau Z)^{-} Z^\tau X$, where
% $(Z^\tau Z)^{-}$ is any generalized inverse of $Z^\tau Z$.
%
% Assumption A1: $\varepsilon \sim N_n(0, \sigma^2 I_n)$ with unknown $\sigma^2 > 0$.
% Assumption A3: $E(\varepsilon) = 0$, $\mathrm{Var}(\varepsilon)$ is an unknown matrix.
%
% $\mathcal{R}(A)$ denotes the smallest linear subspace containing all rows of $A$.

\begin{theorem}[Estimability of $l^\tau \beta$ in linear regression]
\label{thm:3.6}
Assume model (3.25) with assumption A3.
\begin{enumerate}
\item[(i)] A necessary and sufficient condition for $l \in \mathcal{R}^p$ being
  $Q^\tau c$ for some $c \in \mathcal{R}^r$ is
  $l \in \mathcal{R}(Z) = \mathcal{R}(Z^\tau Z)$,
  where $Q$ is given by (3.30) and $\mathcal{R}(A)$ is the smallest linear
  subspace containing all rows of $A$.
\item[(ii)] If $l \in \mathcal{R}(Z)$, then the LSE $l^\tau \hat{\beta}$ is
  unique and unbiased for $l^\tau \beta$.
\item[(iii)] If $l \notin \mathcal{R}(Z)$ and assumption A1 holds, then
  $l^\tau \beta$ is not estimable.
\end{enumerate}
\end{theorem}

\begin{proof}
(i) Note that $a \in \mathcal{R}(A)$ iff $a = A^\tau b$ for some vector $b$.
If $l = Q^\tau c$, then
\[
  l = Q^\tau c
    = Q^\tau Z_*^\tau Z_* (Z_*^\tau Z_*)^{-1} c
    = Z^\tau \bigl[ Z_* (Z_*^\tau Z_*)^{-1} c \bigr],
\]
hence $l \in \mathcal{R}(Z)$. Conversely, if $l \in \mathcal{R}(Z)$, then
$l = Z^\tau \zeta$ for some $\zeta$ and
\[
  l = (Z_* Q)^\tau \zeta = Q^\tau c, \qquad \text{with } c = Z_*^\tau \zeta.
\]

(ii) If $l \in \mathcal{R}(Z) = \mathcal{R}(Z^\tau Z)$, then $l = Z^\tau Z \zeta$
for some $\zeta$, and by (3.29),
\[
  E(l^\tau \hat{\beta})
    = E\bigl[ l^\tau (Z^\tau Z)^{-} Z^\tau X \bigr]
    = \zeta^\tau Z^\tau Z (Z^\tau Z)^{-} Z^\tau Z \beta
    = \zeta^\tau Z^\tau Z \beta
    = l^\tau \beta.
\]
If $\bar{\beta}$ is any other LSE of $\beta$, then by (3.27),
\[
  l^\tau \hat{\beta} - l^\tau \bar{\beta}
    = \zeta^\tau (Z^\tau Z)(\hat{\beta} - \bar{\beta})
    = \zeta^\tau (Z^\tau X - Z^\tau X) = 0.
\]

(iii) Under assumption A1, if $h(X,Z)$ is unbiased for $l^\tau \beta$, then
\[
  l^\tau \beta = \int_{\mathcal{R}^n} h(x,Z) (2\pi)^{-n/2} \sigma^{-n}
    \exp\!\Bigl\{ -\tfrac{1}{2\sigma^2} \| x - Z\beta \|^2 \Bigr\} \, dx.
\]
Differentiating w.r.t.\ $\beta$ and applying Theorem 2.1 yields
\[
  l = Z^\tau \int_{\mathcal{R}^n} h(x,Z) (2\pi)^{-n/2} \sigma^{-n-2}
    (x - Z\beta) \exp\!\Bigl\{ -\tfrac{1}{2\sigma^2} \| x - Z\beta \|^2 \Bigr\} \, dx,
\]
which implies $l \in \mathcal{R}(Z)$.
\end{proof}

\end{document}
